\chapter{Luas dan Volume}\label{ch:g7:areas}

Kelas 6 mengukur persegi panjang, \omterm{def:g6:shapes:triangles}{segitiga siku-siku} dan \omterm{def:g5:solids:def}{balok}
(\cref{ch:g6:measure}). Dengan sepasang gunting --- potong sepotong di
sini, tempel di sana --- gagasan yang sama kini memberi luas
\omterm{def:g7:central:parallelogram}{jajargenjang}, luas \emph{semua} segitiga dan luas cakram, serta \omterm{def:g6:measure:volume}{volume}
\omterm{def:g7:areas:prism}{prisma} dan \omterm{def:g7:areas:prism}{tabung}.

\section{Luas jajargenjang}

\begin{theorem}[Jajargenjang]\label{thm:g7:areas:parallelogram}
\omterm{def:g7:central:parallelogram}{Jajargenjang} dengan sisi sepanjang $b$ (\emph{alasnya}) dan jarak $h$
antara sisi itu dan sisi yang berhadapan (\emph{tingginya}) punya luas
\[
A = b \times h .
\]
\end{theorem}

\begin{proof}[Bukti dengan gunting]
Potonglah \omterm{def:g6:shapes:triangles}{segitiga siku-siku} yang menonjol di satu sisi lalu tempelkan
di sisi yang lain: \omterm{def:g7:central:parallelogram}{jajargenjangnya} menjadi persegi panjang
$b \times h$ dengan luas yang sama.
\end{proof}

\begin{omfigure}
\begin{tikzpicture}[scale=0.82]
  \fill[omDef!14] (0,0) -- (4,0) -- (5.2,1.8) -- (1.2,1.8) -- cycle;
  \draw[thick] (0,0) -- (4,0) -- (5.2,1.8) -- (1.2,1.8) -- cycle;
  \fill[omThm!25] (0,0) -- (1.2,1.8) -- (1.2,0) -- cycle;
  \draw[omThm, thick] (0,0) -- (1.2,1.8) -- (1.2,0) -- cycle;
  \draw[dashed] (1.2,0) -- (1.2,1.8);
  \node[below, font=\small] at (2.0,0) {$b$};
  \draw[Stealth-Stealth] (6.0,0) -- (6.0,1.8)
    node[midway, right, font=\small] {$h$};
  \begin{scope}[xshift=8.2cm]
  \fill[omDef!14] (0,0) rectangle (4,1.8);
  \draw[thick] (0,0) rectangle (4,1.8);
  \fill[omThm!25] (2.8,0) -- (4,1.8) -- (4,0) -- cycle;
  \draw[omThm, thick] (2.8,0) -- (4,1.8) -- (4,0) -- cycle;
  \node[below, font=\small] at (2.0,0) {$b$};
  \end{scope}
\end{tikzpicture}

{\small Bukti dengan gunting: menggeser segitiga merah dari ujung kiri
ke ujung kanan mengubah \omterm{def:g7:central:parallelogram}{jajargenjangnya} menjadi persegi panjang ---
luasnya sama, $b \times h$.}
\end{omfigure}

\begin{remark}
Tingginya adalah jarak \emph{\omterm{def:g6:lines:perp}{tegak lurus}} antara kedua alasnya, bukan
panjang sisi miringnya! \omterm{def:g7:central:parallelogram}{Jajargenjang} yang sangat miring dapat punya
sisi yang panjang tetapi luas yang kecil.
\end{remark}

\section{Luas segitiga}

\begin{theorem}[Segitiga]\label{thm:g7:areas:triangle}
Segitiga dengan alas $b$ dan tinggi $h$ yang bersesuaian (jarak \omterm{def:g6:lines:perp}{tegak
lurus} dari \omterm{def:g5:solids:def}{titik sudut} yang berhadapan ke garis alasnya) punya
luas
\[
A = \frac{b \times h}{2} .
\]
\end{theorem}

\begin{proof}
Dua salinan segitiga itu, yang satu diputar setengah putaran terhadap
titik tengah salah satu sisinya, cocok menjadi sebuah \omterm{def:g7:central:parallelogram}{jajargenjang}
dengan alas $b$ dan tinggi $h$ (itulah \omterm{def:g7:central:def}{simetri pusat} yang sedang
bekerja, \cref{prop:g7:central:preserved}). Segitiganya adalah
setengah darinya: $\frac{bh}{2}$.
\end{proof}

\begin{omfigure}
\begin{tikzpicture}[scale=0.9]
  \fill[omDef!14] (0,0) -- (4,0) -- (1.4,2.1) -- cycle;
  \draw[thick] (0,0) -- (4,0) -- (1.4,2.1) -- cycle;
  \fill[omThm!12] (4,0) -- (1.4,2.1) -- (5.4,2.1) -- cycle;
  \draw[thick, omThm, dashed] (4,0) -- (5.4,2.1) -- (1.4,2.1);
  \draw[dashed] (1.4,0) -- (1.4,2.1);
  \draw[thin] (1.4,0) ++(-0.2,0) -- ++(0,0.2) -- ++(0.2,0);
  \node[below, font=\small] at (2,0) {$b$};
  \node[right, font=\small] at (1.42,1.0) {$h$};
\end{tikzpicture}

{\small Sebuah segitiga dan salinannya yang terputar setengah putaran
(garis putus-putus) membentuk \omterm{def:g7:central:parallelogram}{jajargenjang}: luas segitiganya
$\frac{b\times h}{2}$. Sisi mana pun dari ketiganya dapat menjadi
alasnya --- dengan tingginya \emph{sendiri}.}
\end{omfigure}

\begin{example}\label{ex:g7:areas:triangle}
Sebuah segitiga punya alas $9$ cm dan tinggi yang bersesuaian
berukuran $4$ cm:
\[
A = \frac{9 \times 4}{2} = \frac{36}{2} = 18 \text{ cm}^2 .
\]
Untuk segitiga \emph{siku-siku}, kedua sisi siku-sikunya adalah alas
dan tingginya: kita menemukan kembali rumus
\cref{prop:g6:measure:formulas}.
\end{example}

\section{Luas cakram}

\begin{theorem}[Cakram]\label{thm:g7:areas:disk}
Cakram berjari-jari $r$ punya luas
\[
A = \pi r^2 .
\]
\end{theorem}

\begin{proof}[Gagasan buktinya]
Potonglah cakramnya menjadi banyak juring tipis yang sama lalu
tatalah berselang-seling: juringnya hampir membentuk \omterm{def:g7:central:parallelogram}{jajargenjang}
setinggi $r$ yang alasnya setengah \omterm{def:g6:measure:perimeter}{kelilingnya}, yaitu $\pi r$
(\cref{prop:g6:measure:circle}). Luasnya dekat dengan
$\pi r \times r = \pi r^2$, dan hampirannya menjadi sempurna ketika
juringnya makin tipis. Bukti yang lengkap perlu hitung integral, dari
buku Sekolah Menengah.
\end{proof}

\begin{omfigure}
\begin{tikzpicture}[scale=0.85]
  \foreach \a in {0,45,...,315}
    \fill[omDef!20, draw=omDef] (0,0) -- (\a:1.4)
      arc (\a:\a+45:1.4) -- cycle;
  \begin{scope}[xshift=4.3cm, yshift=-0.15cm]
  \foreach \i in {0,1,2,3} {
    \fill[omDef!20, draw=omDef]
      ({\i*1.1},0) -- ({\i*1.1+0.55},1.4) -- ({\i*1.1+1.1},0) -- cycle;
    \fill[omDef!20, draw=omDef]
      ({\i*1.1+0.55},1.4) -- ({\i*1.1+1.1},0) -- ({\i*1.1+1.65},1.4)
      -- cycle;
  }
  \draw[Stealth-Stealth] (5.2,0) -- (5.2,1.4)
    node[midway, right, font=\small] {$\approx r$};
  \node[below, font=\small] at (2.5,-0.1) {alas $\approx \pi r$};
  \end{scope}
\end{tikzpicture}

{\small Memotong cakramnya lalu membentangkan juringnya: hampir
menjadi \omterm{def:g7:central:parallelogram}{jajargenjang} dengan alas $\pi r$ (setengah tepinya) dan tinggi $r$.}
\end{omfigure}

\begin{example}\label{ex:g7:areas:disk}
Permukaan meja bundar berjari-jari $60$ cm. Luasnya
$\pi \times 60^2 = 3600\pi \approx 11\,310$ cm$^2$, sedikit lebih dari
$1.1$ m$^2$. Perhatikan bedanya: \emph{\omterm{def:g6:measure:perimeter}{kelilingnya}} $2\pi r$ memakai
$r$, sedangkan \emph{luasnya} $\pi r^2$ memakai $r^2$.
\end{example}

\section{Prisma dan tabung}

\begin{definition}[Prisma, tabung]\label{def:g7:areas:prism}
\emph{Prisma}\index{prisma} adalah \omterm{def:g5:solids:def}{bangun ruang} dengan dua sisi
bersegi banyak yang \omterm{def:g6:lines:perp}{sejajar} dan identik (\emph{alasnya}) yang
dihubungkan persegi panjang; \emph{tabung}\index{tabung} adalah bangun
yang sama dengan cakram sebagai alasnya. \emph{Tingginya} adalah jarak
antara kedua alasnya.
\end{definition}

\begin{theorem}[Volume]\label{thm:g7:areas:volume}
Untuk \omterm{def:g7:areas:prism}{prisma} atau \omterm{def:g7:areas:prism}{tabung} dengan luas alas $B$ dan tinggi $h$:
\[
V = B \times h .
\]
\end{theorem}
\admitted

\begin{example}\label{ex:g7:areas:volume}
Sebuah tenda adalah \omterm{def:g7:areas:prism}{prisma} yang berbaring pada sisinya: alasnya
segitiga dengan alas $2$ m dan tinggi $1.5$ m, dan tendanya sepanjang $3$ m.
\begin{enumerate}
  \item Luas alasnya: $B = \frac{2 \times 1.5}{2} = 1.5$ m$^2$.
  \item \omterm{def:g6:measure:volume}{Volumenya}: $V = B \times h = 1.5 \times 3 = 4.5$ m$^3$.
\end{enumerate}
Sebuah kaleng berjari-jari $4$ cm dan tinggi $11$ cm:
$V = \pi \times 4^2 \times 11 = 176\pi \approx 553$ cm$^3$, kira-kira
setengah liter.
\end{example}

\section{Latihan}

\begin{exercise}[$\star$]\label{exo:g7:areas:1}
Hitunglah luas \omterm{def:g7:central:parallelogram}{jajargenjang} dengan alas $8$ cm dan tinggi
$5$ cm; lalu luas \omterm{def:g7:central:parallelogram}{jajargenjang} lain dengan alas $6.5$ m dan tinggi $4$ m.
\end{exercise}

\begin{exercise}[$\star$]\label{exo:g7:areas:2}
Sebuah \omterm{def:g7:central:parallelogram}{jajargenjang} punya sisi $10$ cm dan $6$ cm, dan tingginya
terhadap alas $10$ cm berukuran $4$ cm. Hitunglah luasnya. Mengapa
jawabannya bukan $60$ cm$^2$?
\end{exercise}

\begin{exercise}[$\star$]\label{exo:g7:areas:3}
Hitunglah luas segitiga dengan alas $12$ cm dan tinggi $7$ cm; luas
\omterm{def:g6:shapes:triangles}{segitiga siku-siku} dengan sisi siku-siku $9$ dan $6$; luas segitiga
dengan alas $5.5$ m dan tinggi $2$ m.
\end{exercise}

\begin{exercise}[$\star$]\label{exo:g7:areas:4}
Hitunglah luas cakram berjari-jari $5$ cm, lalu luas setengah cakram
berjari-jari $10$ cm (nilai persis dengan $\pi$, lalu \omterm{def:g4:numbers:round}{dibulatkan} ke satuan).
\end{exercise}

\begin{exercise}[$\star$]\label{exo:g7:areas:5}
Hitunglah \omterm{def:g6:measure:perimeter}{keliling} \emph{dan} luas cakram berjari-jari $3$ cm.
Jawaban yang mana dalam cm dan yang mana dalam cm$^2$?
\end{exercise}

\begin{exercise}[$\star$]\label{exo:g7:areas:6}
Hitunglah \omterm{def:g6:measure:volume}{volume} \omterm{def:g7:areas:prism}{prisma} dengan luas alas $24$ cm$^2$ dan tinggi
$10$ cm; lalu \omterm{def:g6:measure:volume}{volume} \omterm{def:g7:areas:prism}{tabung} berjari-jari $3$ cm dan tinggi $8$ cm
(persis, lalu \omterm{def:g4:numbers:round}{dibulatkan}).
\end{exercise}

\begin{exercise}[$\star$]\label{exo:g7:areas:7}
Sebuah segitiga punya luas $28$ cm$^2$ dan alas $8$ cm. Carilah tinggi
yang bersesuaian, dengan menulis persamaannya lebih dahulu.
\end{exercise}

\begin{exercise}[$\star\star$]\label{exo:g7:areas:8}
Gambarlah segitiga apa saja lalu ukurlah dengan teliti ketiga pasangan
alas--tingginya. Hitunglah $\frac{b \times h}{2}$ untuk tiap
pasangannya: ketiga hasilnya seharusnya sesuai (sampai galat
pengukuran). Mengapa?
\end{exercise}

\begin{exercise}[$\star\star$]\label{exo:g7:areas:9}
Kebun berbentuk trapesium dapat dibelah, oleh satu diagonal, menjadi
dua segitiga yang tingginya sama $h = 20$ m, dengan alas $30$ m dan
$18$ m. Hitunglah luas seluruhnya. Dapatkah kamu menebak rumus umum
untuk trapesium?
\end{exercise}

\begin{exercise}[$\star\star$]\label{exo:g7:areas:10}
Vas berbentuk \omterm{def:g7:areas:prism}{tabung} berjari-jari $6$ cm berisi air dengan tinggi
$15$ cm. Seluruh airnya dituang ke dalam \omterm{def:g5:solids:def}{balok} kosong yang alasnya
$18$ cm $\times$ $12$ cm. Tinggi air berapa yang tercapai di dalam
\omterm{def:g5:solids:def}{baloknya}? (\omterm{def:g6:measure:volume}{Volume} sama, alas berbeda.)
\end{exercise}

\begin{exercise}[$\star\star$]\label{exo:g7:areas:11}
Piza berdiameter $30$ cm berharga $9$ euro; piza berdiameter $40$ cm
berharga $14$ euro. Hitunglah kedua luasnya dan harga per $100$
cm$^2$. Piza mana yang lebih menguntungkan?
\end{exercise}

\begin{exercise}[$\star\star\star$]\label{exo:g7:areas:12}
Kedua sisi siku-siku sebuah \omterm{def:g6:shapes:triangles}{segitiga siku-siku} berukuran $6$ cm dan
$8$ cm, dan sisi miringnya berukuran $10$ cm. Hitunglah luasnya dengan
sisi siku-sikunya, lalu pakailah luas itu untuk mencari tinggi
terhadap \emph{sisi miringnya}.
\end{exercise}

\section{Soal: Persegi yang lenyap}

\begin{problem}[{Soal akhir pekan --- ``bukti'' terkenal bahwa
$64 = 65$, yang dibongkar oleh luas dan kemiringan, dengan piza dan
cincin sebagai pencuci mulut}]
\label{pb:g7:areas:1}
Seorang pesulap memotong persegi $8 \times 8$ menjadi empat
kepingan, menggeser-gesernya, lalu menyusunnya kembali menjadi
persegi panjang $5 \times 13$. Penontonnya membilang:
$8 \times 8 = 64$, tetapi
$5 \times 13 = 65$ --- satu persegi satuan muncul begitu saja
dari udara kosong! Rumus luas bab ini persis alat yang
menangkap triknya. Soal ini lalu mempekerjakan penalaran luas
yang jujur: pada trapesium, pada tumpukan miring, pada
perhitungan piza --- dan akhirnya kembali kepada pesulapnya untuk
pertunjukan kedua yang bahkan lebih aneh.

\medskip
\noindent\textbf{Bagian I --- Triknya, dipentaskan dan dibongkar.}
Keempat kepingan persegi $8 \times 8$ itu adalah: dua \omterm{def:g6:shapes:triangles}{segitiga
siku-siku} dengan sisi siku-siku $8$ dan $3$, dan dua trapesium
siku-siku dengan sisi \omterm{def:g6:lines:perp}{sejajar} $5$ dan $3$ serta tinggi $5$.
\begin{enumerate}
  \item Hitunglah luas persegi aslinya dan luas persegi panjang
        yang dinyatakannya. Nyatakan keanehannya dalam satu baris.
  \item Hitunglah luas tiap kepingan dari keempatnya
        (\cref{thm:g7:areas:triangle}; untuk trapesiumnya,
        potonglah atau pakailah \cref{exo:g7:areas:9}).
  \item Jumlahkan keempat luasnya. Bangun yang mana dari keduanya
        --- persegi atau persegi panjang --- yang sungguh dapat
        diisi kepingannya?
  \item Pada susunan persegi panjangnya, sisi miring sebuah
        segitiga (naik $3$ sepanjang $8$) seharusnya melanjutkan
        tepi miring sebuah trapesium (naik $2$ sepanjang $5$)
        dalam satu ``diagonal'' yang lurus. Ujilah kesegarisannya:
        apakah naik $3$ sepanjang $8$ dan naik $2$ sepanjang $5$
        punya kecuraman yang sama (\cref{thm:g7:prop:cross})?
  \item Jelaskan di mana persegi keenam puluh lima itu bersembunyi:
        apa bentuk celah di sepanjang diagonal palsunya, berapa
        luasnya yang persis, dan mengapa mata melewatkannya?
        (Berapa lebar rata-rata bilah setipis itu kalau
        dibentangkan sepanjang diagonal persegi panjangnya, yang
        kira-kira $13$ satuan panjangnya?)
\end{enumerate}

\medskip
\noindent\textbf{Bagian II --- Penalaran luas yang jujur.}
\begin{enumerate}[resume]
  \item Gambarlah dua \omterm{def:g6:lines:perp}{garis sejajar} yang berjarak $4$ cm, sebuah
        alas sepanjang $6$ cm pada salah satunya, lalu tiga
        segitiga berbeda pada alas itu dengan puncak di berbagai
        titik pada garis yang lain. Hitunglah ketiga luasnya.
        Mengapa semuanya sama, ke mana pun puncaknya bergeser di
        sepanjang garisnya?
  \item Buktikan rumus trapesium secara umum: trapesium dengan
        sisi \omterm{def:g6:lines:perp}{sejajar} $B$ dan $b$ serta tinggi $h$, yang dipotong
        sepanjang sebuah diagonal, memberi dua segitiga dengan
        tinggi yang sama $h$. Simpulkan:
        \[
        A = \frac{(B + b) \times h}{2} .
        \]
  \item Terapkan pada trapesium dengan $B = 9$ cm, $b = 5$ cm,
        $h = 4$ cm.
  \item Tukang kayu memakai rumus yang sama dalam samaran: luas
        $=$ (rata-rata kedua sisi \omterm{def:g6:lines:perp}{sejajarnya}) $\times$ tinggi.
        Jelaskan mengapa itu aturan yang sama, lalu hitunglah
        ulang pertanyaan 8 dengan cara itu.
  \item Miringkan tumpukan kartu remi yang rapi menjadi tumpukan
        yang condong. Jelaskan mengapa \omterm{def:g6:measure:volume}{volumenya} tidak berubah,
        dan apa artinya itu bagi tinggi pada rumus \omterm{def:g7:areas:prism}{prisma}
        $V = B \times h$ (\cref{thm:g7:areas:volume}): tinggi yang
        mana yang harus diambil untuk tumpukan yang condong ---
        panjang condongnya atau panjang tegaknya? (Bandingkan
        \cref{exo:g7:areas:2}, satu dimensi lebih tinggi.)
\end{enumerate}

\medskip
\noindent\textbf{Bagian III --- Cincin, piza, dan kembalinya
pesulap.}
\begin{enumerate}[resume]
  \item Kolam berbentuk lingkaran berjari-jari $3$ m terletak di
        pusat halaman rumput berbentuk lingkaran berjari-jari
        $5$ m. Hitunglah luas cincin rumputnya (persis dengan
        $\pi$, lalu \omterm{def:g4:numbers:round}{dibulatkan} ke m$^2$).
  \item Carilah \omterm{def:g3:shapes:shapes}{jari-jari} cakram tunggal yang luasnya persis sama
        dengan luas cincin itu. (Ketiga \omterm{def:g3:shapes:shapes}{jari-jari} yang kini kamu
        punya membentuk tripel yang terkenal, yang kisahnya
        diceritakan pada \cref{ch:g8:pythagoras}.)
  \item Mana yang lebih banyak piza: satu piza berdiameter
        $40$ cm, atau dua piza berdiameter $28$ cm? Hitunglah lalu
        bandingkan.
  \item Untuk \emph{melipatduakan} luas sebuah piza, kira-kira
        dengan faktor berapa \omterm{def:g3:shapes:shapes}{jari-jarinya} harus tumbuh? Periksalah
        calon $1.4$ (ingatlah bahwa luas mengikuti \emph{kuadrat}
        faktor penskalaannya, \cref{pb:g6:propdata:1}), lalu
        katakan di mana faktor yang persis --- bilangan yang
        kuadratnya $2$ --- membuat penampilan resminya
        (\cref{ex:g8:pythagoras:diagonal}).
  \item Pesulapnya kembali: ia memotong persegi $13 \times 13$
        dengan bilangan $5$, $8$, $13$ memainkan peran yang
        dimainkan $3$, $5$, $8$ sebelumnya, lalu menyusun ulang
        persegi panjang $8 \times 21$. Hitunglah kedua luasnya:
        apa yang lenyap kali ini? Bilangan $3, 5, 8, 13, 21$
        adalah suku berurutan barisan pada
        \cref{pb:g7:fractions:1} --- nyatakan polanya pada
        untung dan rugi pesulapnya, dan pesan moralnya: mengapa
        luas tidak pernah berbohong?

\end{enumerate}
\end{problem}
